% !PRINTED 252
\weberpage{252}

\begin{center}
\S.\,79.

\gesperrt{Die Bezoutiante und ihre Bedeutung für die Wurzelrealität.}
\end{center}

Für die Untersuchung der Realität der Wurzeln einer beliebigen reellen Gleichung in allgemeineren Fällen kann die Function mit Nutzen angewandt werden, die wir im \S.\,73 als Bezoutiante der Gleichung bezeichnet haben.

Ist
\weq{1}{f(x) = x^n + a_1 x^{n-1} + a_2 x^{n-2} + \cdots + a_n = 0}
irgend eine reelle Gleichung $n^{\mathrm{ten}}$ Grades, so war die Bezoutiante folgendermaassen definirt:

Man setze
\weq{2}{y = t_{n-1} f_0(x) + t_{n-2} f_1(x) + \cdots + t_0 f_{n-1}(x),}
worin die $t_0, t_1 \ldots t_{n-1}$ unbestimmte Variable bedeuten und
\[
f_0(x) = 1, \quad f_1(x) = x + a_1, \quad f_2(x) = x^2 + a_1 x + a_2 \ldots
\]
Es ist dann [\S.\,73, (10)]
\weq{3}{S(y^2) = \frac{1}{n}\left[S(y)\right]^2 + B,}
wo $B$ eine quadratische Form der $n - 1$ Variablen $t_0, t_1 \ldots t_{n-2}$ ist mit Coëfficienten, die rational aus den Coëfficienten $a_i$ zusammengesetzt sind. Diese Function $B$ haben wir als Bezoutiante definirt und für die Fälle $n = 4$ und $n = 5$ wirklich gebildet.

Nun sind in der Summe auf der linken Seite von (3)
\weq{4}{y_1^2 + y_2^2 + y_3^2 + \cdots + y_n^2,}
die $y_1, y_2 \ldots y_n$ lineare Functionen der Variablen $t_0, t_1 \ldots t_{n-1}$, und wir haben also in der Formel (3) eine Transformation der quadratischen Form
\weq{5}{\frac{1}{n}\left[S(y)\right]^2 + B = \Phi(t_0, t_1 \ldots t_{n-1})}
% [corrected 2026-08-07, 4x zoom strip 4/4] "auf eine Summe von $n$ Quadraten." is letterspaced --
% the tail of the statement that (5) transforms the quadratic form. Found while calibrating the
% colon check below, not by either audit lane.
\gesperrt{auf eine Summe von $n$ Quadraten.}

Wir machen fürs erste keinerlei beschränkende Voraussetzungen über Gleichheit oder Verschiedenheit der Wurzeln von $f(x)$ und müssen nun zunächst untersuchen, inwieweit die linearen Functionen $y_1, y_2 \ldots y_n$ von einander unabhängig sind. In dieser Beziehung gilt der Satz:
% [verified 2026-08-07, 4x zoom strip 4/4] The page really does end here, on the colon; the Satz
% it announces is printed on 253. So this is a missing continuation marker, NOT lost content --
% the distinction the colon check cannot make on its own.
%CONT>>
